\section{Metodología}
La metodología para esta implementación es Propositiva-Técnica, basada en los hallazgos de los estudios relacionados y el estado actual del SCMC.

\begin{description}
    \item[Fase de Diagnóstico (DMAIC - Definir):] Se define el problema (Articulo 10) como la latencia en la conciliación manual. Se mide el proceso actual (documentado en los requisitos del SCMC) que depende de pagos presenciales en bancos.
    
    \item[Fase de Diseño de Solución (DMAIC - Mejorar):] Se diseña la arquitectura S2S (Servidor-a-Servidor) con Webhooks como la solución técnica. Esto se alinea con el trabajo ya realizado en el SCMC (rutas y migraciones), que prepara al backend C\# para recibir estas notificaciones.
    
    \item[Fase de Justificación (Análisis Financiero):] Se utiliza el modelo financiero del Articulo 10 (VPN, TIR, B/C) para justificar la inversión en horas de desarrollo (C\#/Angular) y licenciamiento de Wompi, demostrando un ROI positivo a través del ahorro en costos operativos.
    
    \item[Definición de Requerimientos No Funcionales:] Se establecen los requisitos de UX/UI (RWD, ícono PSE) basados en la evidencia de Articulo 16, y los requisitos de seguridad (PCI DSS) basados en Articulo 20.
\end{description}   